% Generated draft: requires mathematical and source review.
\documentclass[11pt]{article}
\usepackage[margin=25mm]{geometry}
\usepackage{fontspec}
\setmainfont{Latin Modern Roman}
\usepackage{amsmath,amssymb,mathtools,unicode-math}
\setmathfont{Latin Modern Math}
\usepackage{enumitem,needspace,xurl,hyperref}
\setlength{\parindent}{0pt}
\setlength{\parskip}{5pt}
\setlength{\emergencystretch}{3em}
\newcommand{\N}{\mathbb{N}}
\newcommand{\Z}{\mathbb{Z}}
\newcommand{\Q}{\mathbb{Q}}
\newcommand{\R}{\mathbb{R}}
\newcommand{\C}{\mathbb{C}}
\newcommand{\F}{\mathbb{F}}
\newcommand{\A}{\mathbb{A}}
\newcommand{\PP}{\mathbb{P}}
\newcommand{\E}{\mathbb{E}}
\newcommand{\eps}{\varepsilon}
\newcommand{\dd}{\,\mathrm{d}}
\begin{document}
\section*{UnsolvedMath ID: 10600055}
\textbf{Source problem code:} catalogue entry
\subsection{Definitions and mathematical statement}
A graph is realizable as a chord-intersection graph when it can be obtained from chords on a circle, with two vertices adjacent exactly when the endpoints of the corresponding chords alternate around the circle. This is a usual meaning of realizability in the graph-link setting, but the source does not specify which version of Bouchet's criterion it intends.

Graph-links are equivalence classes of graph diagrams under specified graph Reidemeister-type moves, with any vertex labels or framings required by the source. The record gives neither those moves nor their label conventions. A parity would assign labels, commonly in \(\Z/2\Z\), to graph vertices and satisfy the corresponding move-compatibility conditions. Those conditions cannot be inferred merely from realizability of the underlying unlabelled graph.

Problem-list item 55 asks to construct a parity on graph-links using Bouchet's realizability criterion. It further asks for a parity ``responsible for the cyclically 6-edge connection.'' In a common graph-theoretic convention a graph is cyclically \(6\)-edge-connected when no edge cut of size less than \(6\) separates it into pieces each containing a cycle. This is only a conditional explanation: the source does not define the phrase, identify which associated graph it concerns, or specify what responsibility of a parity means. No equivalence, obstruction theorem, or nontriviality condition is supplied. A precise construction must recover the graph-link move system, Bouchet criterion and intended connection rather than substitute an arbitrary realizability label as a parity.
\subsection{Short English statement}
Build a graph-link parity from the intended Bouchet graph-realizability criterion, and relate it to the source's cyclic \(6\)-edge-connectivity phenomenon. The required move compatibility and the meaning of that relation are missing, along with the in-statement reference.
\subsection{Sources}
\begin{itemize}
\item Local UnsolvedMath registry, permanent ID \(10600055\), exact source code \texttt{catalogue entry}, virtual-knot problem-list item 55. The statement names Bouchet's criterion but has an empty citation after ``see.'' No individual proposer or proposed year is supplied.
\item \emph{Unsolved Problems in Virtual Knot Theory and Combinatorial Knot Theory}, revised survey, dedicated problem-list item 55: \url{https://arxiv.org/abs/1409.2823}. Background: 60 top-level items, compound items retained intact; source-TeX extraction, accessed 2026-07-29.
\item Supplied research summary mentions Bouchet's interlacement/realizability framework and interlace polynomials, with incomplete identifiers arXiv:0209.xxxx (2002) and arXiv:0606.xxxx (2006). No titles or complete references are provided, and the identifiers are not completed here.
\item Source evidence NEEDS\_REVIEW presents an open problem and requires a current primary-literature audit. Supplied OPEN-TRIAGE assessment reports no literature found on this particular application and no resolution; that is a supplied search assessment, not independently verified absence.
\item Dataset attribution: Ulam AI / UnsolvedMath Contributors, CC BY 4.0. Source rights NEEDS\_REVIEW: publicly accessible author/publisher copy, redistribution terms requiring review.
\end{itemize}
\end{document}
